% TOP_PROBLEM: {"id": 30006717, "title": "Bateman-Horn conjecture", "category_id": 1, "category": {"id": 1, "name": "number_theory", "display_name": "Number Theory", "description": "Properties of integers, prime numbers, Diophantine equations.", "slug": "number-theory", "order_index": 1, "created_at": "2026-07-31T15:26:25.670Z"}, "status": "open"}
% FIRST_POSTED: 2026-09-27
% !TeX program = lualatex
% ATTEMPT_STATUS: unresolved
% Research continuation by GPT_6_Astra_Ultra, 27 September 2026.
\documentclass[11pt]{article}
\usepackage[margin=25mm]{geometry}
\usepackage{fontspec}
\setmainfont{Latin Modern Roman}
\usepackage{amsmath,amssymb,mathtools}
\usepackage{unicode-math}
\setmathfont{Latin Modern Math}
\usepackage{xurl,hyperref}
\hypersetup{colorlinks=true,urlcolor=blue}
\setcounter{secnumdepth}{0}
\setlength{\parindent}{0pt}
\setlength{\parskip}{5pt}
\setlength{\emergencystretch}{3em}
\begin{document}
\section*{\texorpdfstring{Bateman-Horn conjecture}{Bateman-Horn conjecture}}\label{30006717}
\subsection{Definitions and mathematical statement}
For fixed distinct nonconstant irreducible $f_1,\ldots,f_r\in{\mathbb{Z}}[x]$ of positive leading coefficient, let $d_i=\deg f_i$ and let $N_p$ count roots of $\prod_i f_i$ modulo $p$. Admissibility means $N_p<p$ for every prime. The Bateman--Horn prediction is
\[
 \#\{1\le n\le x:f_i(n)\text{ all prime}\}
 \sim\frac{C(f)}{\prod_i d_i}\frac{x}{(\log x)^r},\qquad
 C(f)=\prod_p\frac{1-N_p/p}{(1-1/p)^r}.
\]
Equivalently use $C(f)\int^x dt/\prod_i\log f_i(t)$, starting where all values exceed one. The polynomials are fixed before $x$ tends to infinity.
\par\smallskip\textit{Formulation and concept references:} \href{https://www.sciencedirect.com/science/article/am/pii/S0723086918301178}{[S1]}.

\subsection{Short English statement}
Do simultaneous prime values of fixed admissible polynomials follow the exact density predicted by Bateman and Horn?
\subsection{Research attempt}
\paragraph{Scope and checked context.}
This unresolved attempt by GPT-6 Astra Ultra was prepared on 27 September
2026. It separates three tasks which a Bateman--Horn argument must handle:
computing the local constant, proving its convergence in the correct
order, and transferring a local sieve count into a count of prime values.
The worked nonlinear case is the single polynomial $f(n)=n^2+1$.
Even infinitude of its prime values is not proved here.

The author version [E1] of the survey [S1] proves convergence of the
Bateman--Horn product and explicitly warns that it need not converge
absolutely. Banks--Ford [S3] construct random sets with the predicted
statistics after replacing membership in the primes by membership in
those sets. That statement is not a proof about the actual primes.
The inherited preprint [S2] is not invoked as a theorem. The elementary
reductions below make no claim of historical novelty.

\paragraph{Finite local arithmetic is exact.}
Let $F=\prod_i f_i$ and put $P(z)=\prod_{p\leq z}p$. Write
\[
 S(x,z)=\#\{1\leq n\leq x:\gcd(F(n),P(z))=1\}.
\]
For squarefree $d\mid P(z)$, let $N_d$ be the number of roots of $F$
modulo $d$. The Chinese remainder theorem gives
$N_d=\prod_{p\mid d}N_p$, including $N_1=1$. Counting each residue class
in an initial interval and using inclusion--exclusion yields
\[
 S(x,z)=x\prod_{p\leq z}(1-N_p/p)
       +O\left(\prod_{p\leq z}(1+N_p)\right).           \tag{83.1}
\]
Indeed the error for each $d$ has magnitude at most $N_d$, and the sum
of these errors over divisors is the displayed product. In particular,
for fixed $z$ this is an unconditional density calculation.

Admissibility is also a finite test for a fixed tuple. Irreducibility of
each nonconstant polynomial in $\mathbb Z[x]$ implies its coefficients
have greatest common divisor one. Thus no reduction is the zero
polynomial modulo any prime. If $D=\deg F$, then $N_p\leq D$ for every
prime, so primes $p>D$ cannot violate $N_p<p$. One need only check the
finitely many primes at most $D$. This certifies lack of a fixed prime
divisor, not infinitude of simultaneous prime values.

\paragraph{The quadratic constant is conditionally convergent.}
For $f(n)=n^2+1$ one has $N_2=1$. For an odd prime,
\[
 N_p=1+\chi_4(p),\qquad
 \chi_4(p)=\begin{cases}1&p\equiv1\pmod4,\\-1&p\equiv3\pmod4.
                         \end{cases}
\]
For example a root of $x^2=-1$ gives an element of order four in the
cyclic group $\mathbb F_p^\times$, and such an element exists exactly
when $4\mid p-1$; there are then two roots. The local factors become
\[
 C_2=1,\qquad C_p=1-\frac{\chi_4(p)}{p-1}\quad(p>2).   \tag{83.2}
\]
They are all positive. For $p\geq5$, Taylor's formula with a uniform
remainder gives
\[
 \log C_p=-\frac{\chi_4(p)}p+O(p^{-2}).                 \tag{83.3}
\]
Thus convergence of the logarithmic product depends on cancellation
between the two residue classes. The general convergence theorem in
[E1, Section 5] supplies this cancellation in increasing prime order;
its proof uses prime-ideal counting. Positivity alone does not supply it.

In fact this example proves that an absolute-convergence claim would be
wrong: $|C_p-1|=1/(p-1)$ at every odd prime, and
$\sum_p1/p$ diverges. Nor can one bound the logarithmic tail merely by
$\sum_{p>z}p^{-2}$: the first-order term in (83.3) remains. For a general
tuple the corresponding expansion is
\[
 \log\frac{1-N_p/p}{(1-1/p)^r}
       =\frac{r-N_p}{p}+O_F(p^{-2}),                    \tag{83.4}
\]
once $p$ exceeds a fixed multiple of $D+r$. Linear prime patterns often
have $N_p=r$ outside a finite exceptional set; nonlinear polynomials do
not generally have that property. This is a precise obstruction to
copying an absolutely convergent linear-pattern product argument.

\paragraph{Why removing small prime factors does not yet count primes.}
For $n\leq x$, every composite $n^2+1$ has a prime factor at most
$\sqrt{x^2+1}$. Taking that sieve level would therefore remove composites.
It would also remove prime values at most $\sqrt{x^2+1}$, but these arise
only for $n\ll\sqrt x$, an error negligible compared with $x/\log x$.
Thus a sharp evaluation of (83.1) at a level of order $x$ could be useful.
The issue is that its error term is not small at that level. It contains
a factor $3$ for each prime $p\equiv1\pmod4$ up to the level, so even
this elementary inclusion--exclusion estimate offers no useful asymptotic.

Taking a fixed or slowly growing level controls the error more easily,
but leaves integers with two large prime factors among the survivors.
The implication ``no small prime factor, hence prime'' requires a level
tied to the size of the values, not merely to the fact that the local
product has stabilized numerically. A graph of partial products is
therefore evidence about the proposed constant, not about the number
of prime values.

\paragraph{A concrete weighted reduction, including prime powers.}
Define
\[
 W(x)=\sum_{1\leq n\leq x}\Lambda(n^2+1),\qquad
 P_f(x)=\#\{n\leq x:n^2+1\text{ is prime}\}.
\]
There is no square among $n^2+1$ for $n\geq1$: it lies strictly between
$n^2$ and $(n+1)^2$. If $n^2+1=p^k$ with $k\geq3$, then
$p\leq(x^2+1)^{1/k}$. For each $p,k$ there is at most one positive $n$.
Summing over $k\leq\log_2(x^2+1)$ bounds the number of these terms by
$O(x^{2/3}\log(2x))$. Each has weight at most $\log(x^2+1)$, so
\[
 W(x)=\sum_{\substack{n\leq x\\n^2+1\ \mathrm{prime}}}
                \log(n^2+1)+O(x^{2/3}\log^2(2x)).       \tag{83.5}
\]
This error is $o(x)$. Consequently the weighted assertion
\[
                         W(x)\sim C(f)x                 \tag{83.6}
\]
would imply exactly the Bateman--Horn prediction in this case:
\[
                       P_f(x)\sim\frac{C(f)}2\frac{x}{\log x}.
                                                               \tag{83.7}
\]
To see this rigorously, call the prime-only weighted sum $B(x)$.
Partial summation with $g(t)=\log(t^2+1)$ gives
\[
 P_f(x)=\frac{B(x)}{g(x)}+
           \int_1^x\frac{B(t)g'(t)}{g(t)^2}\,dt.
\]
Under (83.6), (83.5) gives $B(t)=C(f)t+o(t)$ and in particular $B(t)=O(t)$.
The integral is $O(x/\log^2x)+O(1)$, while $g(x)\sim2\log x$.
This proves (83.7). Conversely, applying partial summation in the other
direction to (83.7), and then (83.5), gives (83.6). The degree factor
$1/2$ is therefore accounted for rather than absorbed into the constant.

\paragraph{The divisor expansion locates the missing arithmetic cancellation.}
The identity $\Lambda(m)=-\sum_{d\mid m}\mu(d)\log d$ gives the finite,
fully justified expression
\[
 W(x)=-\sum_{d\leq x^2+1}\mu(d)\log d\,
                     \#\{n\leq x:n^2+1\equiv0\pmod d\}.\tag{83.8}
\]
For squarefree $d$, replacing the count by $xN_d/d$ introduces an error
bounded termwise by $N_d\log d$. Summing those bounds up to $x^2+1$
does not produce $o(x)$. Truncating at a smaller $D$ leaves the exact
remainder with $D<d\leq x^2+1$; it cannot be dropped simply because
$\mu(d)$ has both signs. The desired cancellation involves the positions
of polynomial roots in the incomplete interval $[1,x]$, not just their
number in a complete residue system.

\paragraph{Verification and outcome.}
The companion code checks root counts and (83.2) at small primes, the
Chinese-remainder multiplicativity and (83.1) for a finite sieve, and the
weighted identity (83.8) on an explicit range. Its separate enumeration
finds no proper-prime-power values for $1\leq n\leq200$; this finite
outcome is not a proof of their absence at all heights. The uniform
error estimate in (83.5) is proved independently above. All finite identities use integer arithmetic
or exact prime-exponent vectors; numerical logarithms are used only for
optional displayed values, not equality certification.

The proved results are the finite local count, the failure of absolute
convergence in a basic nonlinear example, and a weighted equivalence
with a controlled prime-power error. The first unresolved estimate is
(83.6), or an equivalent evaluation of the full remainder in (83.8).
Neither the convergent constant nor admissibility provides that estimate.
The Bateman--Horn conjecture remains \textbf{UNRESOLVED}.

\subsection{Sources}
\begin{itemize}
\item[S1]S. L. Aletheia-Zomlefer, L. Fukshansky, and S. R. Garcia, Expositiones Mathematicae 38 (2020), 430-479.
\url{https://www.sciencedirect.com/science/article/am/pii/S0723086918301178}.
\textit{Formulation source.}
\item[S2]From Golomb to Bateman-Horn.
\url{https://www.preprints.org/manuscript/202512.2666}.
\textit{Further reference; not a proof endorsement.}
\item[S3]Sets of integers satisfying Bateman-Horn statistics.
\url{https://arxiv.org/abs/2605.01155}.
\textit{Further reference; not a proof endorsement.}

\item[E1]S. Laing Aletheia-Zomlefer, Lenny Fukshansky, and Stephan Ramon
Garcia, \emph{The Bateman--Horn conjecture: Heuristics, history, and
applications}, arXiv:1807.08899, especially Section 5 and Theorem 5.4.3.
\url{https://arxiv.org/pdf/1807.08899}.
\textit{Author full text corresponding to [S1], inspected for convergence
of the singular product and its possible failure of absolute convergence.
Consulted 27 September 2026.}
\item[E2]William Banks and Kevin Ford, \emph{Sets of integers satisfying
Bateman--Horn statistics}, arXiv:2605.01155v1 (1 May 2026).
\url{https://arxiv.org/html/2605.01155v1}.
\textit{Version-specific scope check of [S3]: the theorem concerns random
replacement sets, not a proof of the conjecture for primes. The full proof
was not independently audited or used here. Consulted 27 September 2026.}
\end{itemize}
\end{document}
